% GENERATED FILE. Edit canonical lessons, then run npm run generate:books.
% canonical-source-hash: fnv1a64:52d58475cce8ce17
% canonical-lessons: GE-C315-ritter, GE-C315-koenig, GE-C315-prinzessin, GE-C315-riese, GE-C315-zwerg

\chapter{Ritter und Riesen}
\label{ch:ge-ritter-und-riesen}
\hlchaptermodality{315}

\begin{chapteropening}
\textbf{By the end of this chapter:} \emph{I can name the people of old tales: a knight, a king, a princess, a giant and a dwarf.}

Complete the last lesson of chapter 315: I can name the people of old tales: a knight, a king, a princess, a giant and a dwarf.
\end{chapteropening}

\section[der Ritter]{der Ritter — the knight}

\label{lesson:GE-C315-ritter}



\begin{quote}
Before the new one: say the German for the curse, then the German for the dragon.
\end{quote}

\subsection*{You'll want to know: der Ritter}
\textbf{der Ritter} — \textquotedblleft{}the knight\textquotedblright{}.

\begin{quote}
Der Ritter ritt auf einem weißen Pferd.
\end{quote}

\begin{quote}
\emph{The knight rode a white horse.}
\end{quote}

Plural: \emph{die Ritter}. Built on \emph{reiten}, to ride.

\begin{cousinweb}
\emph{Ritter} comes from Middle High German \emph{ritter}, shaped after Middle Dutch \emph{ridder} (EtymDB), built on the root of \emph{reiten} \textquotedblleft{}to ride\textquotedblright{}: a rider.
\end{cousinweb}

\begin{grammarlens}[title={What you've built}]
One more word for telling, explaining and describing at B1.
\end{grammarlens}

\subsection*{Your turn}
\begin{itemize}
\raggedright
  \item \emph{Say it:} \emph{der Ritter}
  \item \emph{Say it:} \emph{der Ritter}, once more
  \item \emph{Say it:} a sentence of your own with \emph{der Ritter}
  \item \emph{Recall:} say the German for the curse, then the German for the dragon, then say \emph{der Ritter} again
\end{itemize}

\begin{tcolorbox}[breakable,colback=teal!4,colframe=teal!35!black,title={Before you move on}]
What does der Ritter mean? (\textquotedblleft{}The knight\textquotedblright{}.) Say it once more.
\end{tcolorbox}
\section[der König]{der König — the king}

\label{lesson:GE-C315-koenig}



\begin{quote}
Before the new one: say the German for the dragon, then the German for the knight.
\end{quote}

\subsection*{You'll want to know: der König}
\textbf{der König} — \textquotedblleft{}the king\textquotedblright{}.

\begin{quote}
Der König hatte drei Töchter.
\end{quote}

\begin{quote}
\emph{The king had three daughters.}
\end{quote}

Plural: \emph{die Könige}. \emph{Die Königin} — the queen.

\begin{cousinweb}
\emph{König} goes back to Proto-Germanic \emph{kuningaz} \textquotedblleft{}king\textquotedblright{}, which EtymDB builds from \emph{kunją} \textquotedblleft{}kin, family\textquotedblright{} and the ending \emph{-ingaz} \textquotedblleft{}one who comes from\textquotedblright{}: a man of the noble kin.
\end{cousinweb}

\begin{grammarlens}[title={What you've built}]
One more word for telling, explaining and describing at B1.
\end{grammarlens}

\subsection*{Your turn}
\begin{itemize}
\raggedright
  \item \emph{Say it:} \emph{der König}
  \item \emph{Say it:} \emph{der König}, once more
  \item \emph{Say it:} a sentence of your own with \emph{der König}
  \item \emph{Recall:} say the German for the dragon, then the German for the knight, then say \emph{der König} again
\end{itemize}

\begin{tcolorbox}[breakable,colback=teal!4,colframe=teal!35!black,title={Before you move on}]
What does der König mean? (\textquotedblleft{}The king\textquotedblright{}.) Say it once more.
\end{tcolorbox}
\section[die Prinzessin]{die Prinzessin — the princess}

\label{lesson:GE-C315-prinzessin}



\begin{quote}
Before the new one: say the German for the knight, then the German for the king.
\end{quote}

\subsection*{You'll want to know: die Prinzessin}
\textbf{die Prinzessin} — \textquotedblleft{}the princess\textquotedblright{}.

\begin{quote}
Die Prinzessin wollte nicht heiraten.
\end{quote}

\begin{quote}
\emph{The princess did not want to marry.}
\end{quote}

Plural: \emph{die Prinzessinnen}. \emph{Der Prinz} — the prince.

\begin{cousinweb}
\emph{Prinzessin} is built after French \emph{princesse} (EtymDB), from Latin \emph{prīnceps} \textquotedblleft{}first, chief\textquotedblright{}, from \emph{prīmus} \textquotedblleft{}first\textquotedblright{} and \emph{capere} \textquotedblleft{}to take\textquotedblright{}.
\end{cousinweb}

\begin{grammarlens}[title={What you've built}]
One more word for telling, explaining and describing at B1.
\end{grammarlens}

\subsection*{Your turn}
\begin{itemize}
\raggedright
  \item \emph{Say it:} \emph{die Prinzessin}
  \item \emph{Say it:} \emph{die Prinzessin}, once more
  \item \emph{Say it:} a sentence of your own with \emph{die Prinzessin}
  \item \emph{Recall:} say the German for the knight, then the German for the king, then say \emph{die Prinzessin} again
\end{itemize}

\begin{tcolorbox}[breakable,colback=teal!4,colframe=teal!35!black,title={Before you move on}]
What does die Prinzessin mean? (\textquotedblleft{}The princess\textquotedblright{}.) Say it once more.
\end{tcolorbox}
\section[der Riese]{der Riese — the giant}

\label{lesson:GE-C315-riese}



\begin{quote}
Before the new one: say the German for the king, then the German for the princess.
\end{quote}

\subsection*{You'll want to know: der Riese}
\textbf{der Riese} — \textquotedblleft{}the giant\textquotedblright{}.

\begin{quote}
Der Riese war so groß wie ein Haus.
\end{quote}

\begin{quote}
\emph{The giant was as big as a house.}
\end{quote}

Weak noun: \emph{den Riesen}. \emph{Riesig} — huge.

\begin{cousinweb}
\emph{Riese} is inherited from Old High German \emph{riso}, from Proto-Germanic \emph{risiz} \textquotedblleft{}giant\textquotedblright{} (EtymDB).
\end{cousinweb}

\begin{grammarlens}[title={What you've built}]
One more word for telling, explaining and describing at B1.
\end{grammarlens}

\subsection*{Your turn}
\begin{itemize}
\raggedright
  \item \emph{Say it:} \emph{der Riese}
  \item \emph{Say it:} \emph{der Riese}, once more
  \item \emph{Say it:} a sentence of your own with \emph{der Riese}
  \item \emph{Recall:} say the German for the king, then the German for the princess, then say \emph{der Riese} again
\end{itemize}

\begin{tcolorbox}[breakable,colback=teal!4,colframe=teal!35!black,title={Before you move on}]
What does der Riese mean? (\textquotedblleft{}The giant\textquotedblright{}.) Say it once more.
\end{tcolorbox}
\section[der Zwerg]{der Zwerg — the dwarf}

\label{lesson:GE-C315-zwerg}



\begin{quote}
Before the new one: say the German for the princess, then the German for the giant.
\end{quote}

\subsection*{You'll want to know: der Zwerg}
\textbf{der Zwerg} — \textquotedblleft{}the dwarf\textquotedblright{}.

\begin{quote}
Sieben Zwerge wohnten hinter den Bergen.
\end{quote}

\begin{quote}
\emph{Seven dwarfs lived beyond the mountains.}
\end{quote}

Plural: \emph{die Zwerge}.

\begin{cousinweb}
\emph{Zwerg} is inherited from Old High German \emph{twerc}, from Proto-Germanic \emph{dwergaz} (EtymDB), the same source as English \emph{dwarf}.
\end{cousinweb}

\begin{grammarlens}[title={What you've built}]
One more word for telling, explaining and describing at B1.
\end{grammarlens}

\subsection*{Your turn}
\begin{itemize}
\raggedright
  \item \emph{Say it:} \emph{der Zwerg}
  \item \emph{Say it:} \emph{der Zwerg}, once more
  \item \emph{Say it:} a sentence of your own with \emph{der Zwerg}
  \item \emph{Recall:} say the German for the princess, then the German for the giant, then say \emph{der Zwerg} again
\end{itemize}

\begin{tcolorbox}[breakable,colback=teal!4,colframe=teal!35!black,title={Before you move on}]
What does der Zwerg mean? (\textquotedblleft{}The dwarf\textquotedblright{}.) Say it once more.
\end{tcolorbox}
